\documentclass{ctexart}
\usepackage{Note}
\usepackage{unicode-math}
\setCJKmainfont[
  BoldFont = 思源宋体 Bold,
  ItalicFont = 楷体,
]{宋体}
\begin{document}
\section{对称性和守恒律}
\subsection{平移算符}
\begin{definition}[平移算符]
    平移算符$\hat{T}$作用在一个函数上并将其平移一段距离:
    \[\hat{T}(a)\psi(x)=\psi(x-a)\]
\end{definition}
对$\psi(x-a)$进行泰勒展开可得
\[\hat{T}(a)\psi(x)=\psi(x-a)=\sum_{n=0}^{\infty}\dfrac{1}{n!}(-a)^n\dfrac{\d^n\psi(x)}{\d x^n}=\sum_{n=0}^{\infty}\dfrac{1}{n!}\left(-\dfrac{\i a}{\hbar}\hat{p}\right)^n\psi(x)=\exp\left(-\dfrac{\i a}{\hbar}\hat{p}\right)\psi(x)\]
因此有
\[\hat{T}(a)=\exp\left(-\dfrac{\i a}{\hbar}\hat{p}\right)\]
我们说动量是平移算符的\tbf{生成元}.
\end{document}